\subsection{\texorpdfstring{Another interpretation of \(\lambda^{\#}\)}{Another interpretation of \textbackslash lambda\^{}\{\textbackslash\#\}}}

For every \(x \in X\) , the mapping \(\xi \mapsto x\xi\) of H into X is proper (GT, III, §4, No.~2, Prop. 4), therefore \(\beta\) admits an image measure on X under this mapping, which image is concentrated on the orbit xH (Ch. V, §6, No.~2, Cor. 3 of Prop. 2); since \(\beta\) is left-invariant, this image measure depends only on the class \(u = \pi(x)\) of x in X/H, and will be denoted \(\beta_{u}\) . By definition, for \(f \in \mathcal{K}(X)\) ,

\[
\int_ {\mathrm{X}} f (y) d \beta_ {u} (y) = \int_ {\mathrm{H}} f (x \xi) d \beta (\xi) = f ^ {\flat} (u).\tag{8}
\]

We thus see that

\[
(\varepsilon_ {u}) ^ {\sharp} = \beta_ {u}.\tag{9}
\]

Lemma 2. --- Let \(f\) be a function on \(X\) , with values in a topological space.

\begin{enumerate}
\def\labelenumi{\alph{enumi})}
\tightlist
\item
  If \(f\) is a numerical function \(\geqslant 0\) then, for \(x \in X\) ,
\end{enumerate}

\[
\int_ {X} ^ {*} f (y) d \beta_ {\dot {x}} (y) = \int_ {H} ^ {*} f (x \xi) d \beta (\xi) \quad (\dot {x} = \pi (x)).
\]

\begin{enumerate}
\def\labelenumi{\alph{enumi})}
\setcounter{enumi}{1}
\item
  For \(f\) to be \(\beta_{\dot{x}}\) -measurable, it is necessary and sufficient that the function \(\xi \mapsto f(x\xi)\) on \(H\) be \(\beta\) -measurable.
\item
  Suppose that \(\mathbf{f}\) is a function on \(\mathbf{X}\) , with values in a Banach space or in \(\overline{\mathbf{R}}\) ; then, for \(\mathbf{f}\) to be \(\beta_{\dot{x}}\) -integrable (resp. essentially \(\beta_{\dot{x}}\) -integrable), it is necessary and sufficient that the function \(\xi \mapsto \mathbf{f}(x\xi)\) on \(\mathbf{H}\) be \(\beta\) -integrable (resp. essentially \(\beta\) -integrable), in which case \(\int_{\mathbf{X}} \mathbf{f}(y) d\beta_{\dot{x}}(y) = \int_{\mathbf{H}} \mathbf{f}(x\xi) d\beta(\xi)\) .
\end{enumerate}

This follows from Ch. V, §4, Prop. 2, Prop. 3 and Th. 2.

Since \(f^{\flat} \in \mathcal{K}(X/H)\) for \(f \in \mathcal{K}(X)\) , formula (8) proves that the mapping \(u \mapsto \beta_{u}\) of X/H into \(\mathcal{M}(X)\) is vaguely continuous, that the family \((\beta_{u})\) is \(\lambda\) -adequate \(^{1}\) for any positive measure \(\lambda\) on X/H, and that

\[
\lambda^ {\sharp} = \int_ {\mathrm{X/H}} \beta_ {u} d \lambda (u),\tag{10}
\]

which furnishes a new interpretation of \(\lambda^{\sharp}\) .

PROPOSITION 5. --- Let \(\lambda\) be a positive measure on X/H.

\begin{enumerate}
\def\labelenumi{\alph{enumi})}
\item
  Let \(f\) be a \(\lambda^{\sharp}\) -measurable function on \(X\) , with values in a topological space, constant outside a countable union of \(\lambda^{\sharp}\) -integrable sets. Then, the set of \(\dot{x} \in X / H\) such that the function \(\xi \mapsto f(x\xi)\) is not \(\beta\) -measurable is locally \(\lambda\) -negligible.
\item
  Let \(f\) be a \(\lambda^{\sharp}\) -measurable function \(\geqslant 0\) on \(X\) , zero outside a countable union of \(\lambda^{\sharp}\) -integrable sets. Then, the function \(\dot{x} \mapsto \int^{*} f(x\xi) d\beta(\xi)\) on \(X / H\) is \(\lambda\) -measurable, and
\end{enumerate}

\[
\int_ {\mathrm{X}} ^ {*} f (x) d \lambda^ {\sharp} (x) = \int_ {\mathrm{X/H}} ^ {*} d \lambda (\dot {x}) \int_ {\mathrm{H}} ^ {*} f (x \xi) d \beta (\xi) \qquad \big (\dot {x} = \pi (x) \big).
\]

\begin{enumerate}
\def\labelenumi{\alph{enumi})}
\setcounter{enumi}{2}
\tightlist
\item
  Let \(\mathbf{f}\) be a \(\lambda^{\sharp}\) -integrable function on \(X\) , with values in a Banach space or in \(\overline{\mathbf{R}}\) . Then, the set of \(\dot{x} \in X / H\) such that \(\xi \mapsto \mathbf{f}(x\xi)\) is not \(\beta\) -integrable is \(\lambda\) -negligible; the function \(\mathbf{f}^{\flat}\) on \(X / H\) defined almost everywhere by the formula
\end{enumerate}

\[
\mathbf {f} ^ {\flat} (\dot {x}) = \int_ {\mathrm{H}} \mathbf {f} (x \xi) d \beta (\xi) \quad \left(\dot {x} = \pi (x)\right)\tag{11}
\]

is \(\lambda\) -integrable, and

\[
\int_ {\mathrm{X/H}} \mathbf {f} ^ {\flat} d \lambda = \int_ {\mathrm{X}} \mathbf {f} d \lambda^ {\sharp}\tag{12}
\]

and

\[
\int_ {\mathrm{X} / \mathrm{H}} | \mathbf {f} ^ {\flat} | d \lambda \leqslant \int_ {\mathrm{X}} | \mathbf {f} | d \lambda^ {\sharp}.\tag{13}
\]

\begin{enumerate}
\def\labelenumi{\alph{enumi})}
\setcounter{enumi}{3}
\tightlist
\item
  Let \(\mathbf{f}\) be a \(\lambda^{\sharp}\) -measurable function on \(\mathbf{X}\) , with values in a Banach space or in \(\overline{\mathbf{R}}\) , and zero outside a countable union of \(\lambda^{\sharp}\) -integrable sets. Then, for \(\mathbf{f}\) to be \(\lambda^{\sharp}\) -integrable, it is necessary and sufficient that
\end{enumerate}

\[
\int_ {\mathrm{X} / \mathrm{H}} ^ {*} d \lambda (\dot {x}) \int_ {\mathrm{H}} ^ {*} | \mathbf {f} (x \xi) | d \beta (\xi) <   + \infty \quad \left(\dot {x} = \pi (x)\right).
\]

Taking into account Lemma 2, the assertions a), b) and c) follow from Ch. V, §3, Prop. 4, Prop. 5 and Th. 1 (with the exception of (13), which follows from (12) because it is clear that \(|\mathbf{f}^{\flat}| \leqslant |\mathbf{f}|^{\flat}\) ); d) follows from b).

PROPOSITION 6. --- Let \(\lambda\) be a positive measure on X/H.

\begin{enumerate}
\def\labelenumi{\alph{enumi})}
\item
  Let N be a subset of X/H. For N to be locally \(\lambda\) -negligible, it is necessary and sufficient that \(\overline{\pi}^{-1}(N)\) be locally \(\lambda^{\sharp}\) -negligible.
\item
  Let \(g\) be a function on \(\mathrm{X / H}\) , with values in a topological space. For \(g\) to be \(\lambda\) -measurable, it is necessary and sufficient that \(g \circ \pi\) be \(\lambda^{\sharp}\) -measurable.
\item
  Let \(\mathbf{h}\) be a function on \(\mathrm{X / H}\) , with values in a Banach space or in \(\overline{\mathbf{R}}\) . For \(\mathbf{h}\) to be locally \(\lambda\) -integrable, it is necessary and sufficient that \(\mathbf{h} \circ \pi\) be locally \(\lambda^{\sharp}\) -integrable, in which case \((\mathbf{h} \cdot \lambda)^{\sharp} = (\mathbf{h} \circ \pi) \cdot \lambda^{\sharp}\) .
\end{enumerate}

Suppose \(\mathbf{h} \circ \pi\) is locally \(\lambda^{\sharp}\) -integrable. For every \(f \in \mathcal{K}(X)\) , \(f \cdot (\mathbf{h} \circ \pi)\) is \(\lambda^{\sharp}\) -integrable, therefore (Prop. 5) the function \((f \cdot (\mathbf{h} \circ \pi))^b = f^b \cdot \mathbf{h}\) is \(\lambda\) -integrable and

\[
\int_ {\mathrm{X} / \mathrm{H}} f ^ {\flat} \cdot \mathbf {h} d \lambda = \int_ {\mathrm{X}} f \cdot (\mathbf {h} \circ \pi) d \lambda^ {\sharp}.
\]

Since \(f \mapsto f^{\flat}\) is a surjective mapping of \(\mathcal{K}(\mathrm{X})\) onto \(\mathcal{K}(\mathrm{X}/\mathrm{H})\) , this shows that \(\mathbf{h}\) is locally \(\lambda\) -integrable and that

\[
(\mathbf {h} \cdot \lambda) ^ {\sharp} = (\mathbf {h} \circ \pi) \cdot \lambda^ {\sharp}.
\]

In particular, if \(\overline{\pi}^{1}(N)\) is locally \(\lambda^{\sharp}\) -negligible, then \(\varphi_{N} \circ \pi\) is locally \(\lambda^{\sharp}\) -negligible, therefore \((\varphi_{N} \cdot \lambda)^{\sharp} = (\varphi_{N} \circ \pi) \cdot \lambda^{\sharp} = 0\) , consequently \(\varphi_{N} \cdot \lambda = 0\) and \(N\) is locally \(\lambda\) -negligible. Now suppose that \(g \circ \pi\) is \(\lambda^{\sharp}\) -measurable. Let \(K'\) be a compact subset of \(X/H\) . Let \(f \in \mathcal{K}_{+}(X)\) be such that \(f^{\flat} = 1\) on \(K'\) (No.~1, Prop. 2), and let \(K = \operatorname{Supp} f\) ; one has \(\pi(K) \supset K'\) . There exists a partition of \(K\) formed of a \(\lambda^{\sharp}\) -negligible set \(M\) and a sequence \((K_n)\) of compact sets such that \((g \circ \pi)|K_n\) is continuous for every \(n\) . Then \(g|\pi(K_n)\) is continuous. Let \(P\) be the set of points of \(K'\) not belonging to \(\pi(K_1) \cup \pi(K_2) \cup \ldots\) ; then \(\overline{\pi}^{1}(P) \cap K\) is contained in \(M\) , hence is \(\lambda^{\sharp}\) -negligible; therefore \(f \cdot \varphi_{\overline{\pi}(P)}^{-1}\) is \(\lambda^{\sharp}\) -negligible; it follows (Prop. 5) that

\[
0 = \int_ {\mathrm{X}} f \cdot \varphi_ {\pi^ {- 1} (\mathrm{P})} d \lambda^ {\sharp} = \int_ {\mathrm{X/H}} f ^ {\flat} \cdot \varphi_ {\mathrm{P}} d \lambda \geqslant \int_ {\mathrm{X/H}} ^ {*} \varphi_ {\mathrm{P}} d \lambda ,
\]

therefore \(\mathbf{P}\) is \(\lambda\) -negligible, and \(g\) is \(\lambda\) -measurable.

If N is locally \(\lambda\) -negligible, then \(\overline{\pi}^{1}(N)\) is locally \(\lambda^{\sharp}\) -negligible (Ch. V, §6, No.~6, Cor. 1 of Prop. 10). If g is \(\lambda\) -measurable, then \(g \circ \pi\) is \(\lambda^{\sharp}\) -measurable (ibid.). Finally, suppose h is locally \(\lambda\) -integrable. Then we already know that \(h \circ \pi\) is \(\lambda^{\sharp}\) -measurable. For every \(f \in \mathcal{K}_{+}(X)\) we have, by Prop. 5,

\[
\int_ {X} ^ {*} f (x) | \mathbf {h} | (\pi (x)) d \lambda^ {\sharp} (x) = \int_ {X / H} ^ {*} | \mathbf {h} | (u) f ^ {\flat} (u) d \lambda (u) <   + \infty ,
\]

therefore \(\mathbf{h} \circ \pi\) is locally \(\lambda^{\sharp}\) -integrable.

COROLLARY 1. --- Let \(\lambda, \lambda'\) be two positive measures on X/H. For \(\lambda'\) to have base \(\lambda\) , it is necessary and sufficient that \(\lambda'^{\sharp}\) have base \(\lambda^{\sharp}\) . For \(\lambda\) and \(\lambda'\) to be equivalent, it is necessary and sufficient that \(\lambda^{\sharp}\) and \(\lambda'^{\sharp}\) be equivalent.

The first assertion follows from Prop. 6, a) and c). The second follows from the first.

COROLLARY 2. --- Let \(\lambda\) be a positive measure on X/H, and \(f\) a \(\lambda^{\sharp}\) -measurable numerical function on X. Suppose that, for every \(\xi \in H\) , \(\delta(\xi)f = f\) locally \(\lambda^{\sharp}\) -almost everywhere. Then, there exists a \(\lambda\) -measurable function \(g\) on X/H such that \(f = g \circ \pi\) locally \(\lambda^{\sharp}\) -almost everywhere.

Replacing \(f\) by \(f / (1 + |f|)\) , one reduces to the case that \(f\) is bounded, hence locally \(\lambda^{\sharp}\) -integrable. Let \(\mu = f \cdot \lambda^{\sharp}\) . The hypothesis on \(f\) implies that \(\pmb{\delta}(\xi) \mu = f \cdot \pmb{\delta}(\xi) \lambda^{\sharp} = \Delta_{\mathrm{H}}(\xi) \mu\) for all \(\xi \in \mathrm{H}\) . There then exists (Prop. 4) a measure \(\lambda'\) on \(\mathrm{X} / \mathrm{H}\) such that \(\mu = \lambda'^{\sharp}\) . By Cor. 1, there exists a locally \(\lambda\) -integrable function \(g\) on \(\mathrm{X} / \mathrm{H}\) such that \(\lambda' = g \cdot \lambda\) . By Prop. 6, \(f \cdot \lambda^{\sharp} = \lambda'^{\sharp} = (g \circ \pi) \cdot \lambda^{\sharp}\) , whence \(f = g \circ \pi\) locally \(\lambda^{\sharp}\) -almost everywhere.

COROLLARY 3. --- a) Let \((\lambda_{\iota})_{\iota \in I}\) be a family of real measures on X/H. For the family \((\lambda_{\iota})\) to be bounded above in \(\mathcal{M}(\mathrm{X / H})\) , it is necessary and sufficient that the family \((\lambda_{\iota}^{\sharp})\) be bounded above in \(\mathcal{M}(\mathrm{X})\) , in which case

\[
\sup (\lambda_ {\iota} ^ {\sharp}) = (\sup \lambda_ {\iota}) ^ {\sharp}.
\]

\begin{enumerate}
\def\labelenumi{\alph{enumi})}
\setcounter{enumi}{1}
\item
  Let \(\lambda\) be a real measure on X/H. Then \((\lambda^{+})^{\sharp} = (\lambda^{\sharp})^{+}\) and \((\lambda^{-})^{\sharp} = (\lambda^{\sharp})^{-}\) .
\item
  Let \(\lambda\) be a complex measure on X/H. Then \(|\lambda|^{\sharp} = |\lambda^{\sharp}|\) .
\end{enumerate}

Assume the family \((\lambda_{\iota})\) to be bounded above and let \(\mu = \sup \lambda_{\iota}\) . Since \(\lambda \geqslant 0\) implies \(\lambda^{\sharp} \geqslant 0\) , we have \(\mu^{\sharp} \geqslant \lambda_{\iota}^{\sharp}\) for all \(\iota\) , which shows that the family \((\lambda_{\iota}^{\sharp})\) is bounded above and that

\[
(\sup \lambda_ {\iota}) ^ {\sharp} \geqslant \sup (\lambda_ {\iota} ^ {\sharp}).
\]

Conversely, assume the family \((\lambda_{\iota}^{\sharp})\) to be bounded above and let \(\nu = \sup(\lambda_{\iota}^{\sharp})\) . Since \(\delta(\xi)\lambda_{\iota}^{\sharp} = \Delta_{\mathrm{H}}(\xi)\lambda_{\iota}^{\sharp}\) for all \(\xi \in H\) , obviously \(\delta(\xi)\nu = \Delta_{\mathrm{H}}(\xi)\nu\) , therefore there exists a measure \(\mu' \in \mathcal{M}(\mathrm{X}/\mathrm{H})\) such that \(\nu = \mu'^{\sharp}\) . Since \(\lambda^{\sharp} \geqslant 0\) implies \(\lambda \geqslant 0\) , we have \(\mu' \geqslant \lambda_{\iota}\) for all \(\iota\) , which shows that the family \((\lambda_{\iota})\) is bounded above and that \(\nu = \mu'^{\sharp} \geqslant (\sup \lambda_{\iota})^{\sharp}\) , whence

\[
\sup \left(\lambda_ {\iota} ^ {\sharp}\right) \geqslant \left(\sup \lambda_ {\iota}\right) ^ {\sharp},
\]

which completes the proof of a). The assertion b) then follows at once since, for example, \(\lambda^{+}\) is none other than \(\sup(\lambda,0)\) . To prove c), it suffices to note that \(|\lambda| = \sup \mathcal{R}(\alpha \lambda)\) over the complex numbers \(\alpha\) of absolute value 1, and on the other hand that \(\mathcal{R}(\mu^{\sharp}) = (\mathcal{R}\mu)^{\sharp}\) for every \(\mu \in \mathcal{M}(\mathrm{X / H})\) .

Remarks. --- 1) Prop. 6 a) may be expressed by saying that \(\lambda\) is a pseudo-image measure of \(\lambda^{\sharp}\) under \(\pi\) (Ch. VI, §3, No.~2, Def. 1).

\begin{enumerate}
\def\labelenumi{\arabic{enumi})}
\setcounter{enumi}{1}
\tightlist
\item
  Suppose H is compact and \(\beta\) is normalized. The saturation of every compact subset of X is compact. Therefore, if \(f \in \mathcal{K}(\mathrm{X}/\mathrm{H})\) then \(f \circ \pi \in \mathcal{K}(\mathrm{X})\) ; and, for every positive measure \(\lambda\) on X/H, Prop. 5 c) gives
\end{enumerate}

\[
\int_ {\mathrm{X}} (f \circ \pi) (x) d \lambda^ {\sharp} (x) = \int_ {\mathrm{X/H}} f (u) d \lambda (u).
\]

In other words, \(\lambda\) is the image of \(\lambda^{\sharp}\) under \(\pi\) .

\begin{enumerate}
\def\labelenumi{\arabic{enumi})}
\setcounter{enumi}{2}
\item
  Cor. 3 c) of Prop. 6 shows at once that the results of this subsection remain valid in the case of complex measures (except for those that involve the upper integral).
\item
  Let \(\mathbf{m}\) be a vectorial measure on X/H with values in E, and let \(q\) be a lower semi-continuous semi-norm on E. For \(\mathbf{m}\) to be \(q\) -majorizable (Ch. VI, §2, No.~3, Def. 3), it is necessary and sufficient that \(\mathbf{m}^{\sharp}\) be so, in which case \(q(\mathbf{m}^{\sharp}) = q(\mathbf{m})^{\sharp}\) . This follows at once from the definitions and Cor. 3 a).
\end{enumerate}

On the other hand, let \(\mu\) be a positive measure on X/H. For m to be scalarly of base \(\mu\) , it is necessary and sufficient that \(\mathbf{m}^{\sharp}\) be scalarly of base \(\mu^{\sharp}\) : this follows from Cor. 1.

Finally, if \(\mathbf{m}\) has base \(\mu\) , with density \(\mathbf{f}\) with respect to \(\mu\) (Ch. VI, §2, No.~4, Def. 4), then \(\mathbf{m}^{\sharp}\) has base \(\mu^{\sharp}\) , with density \(\mathbf{f} \circ \pi\) : this follows from Prop. 6 c).

